\scanpage{71}
% lines 1–3, opening computations
\[
\begin{aligned}
&\deleted{\rho(\overbrace{u_1\otimes x_2+x_1\otimes u_2}^{\unsure{\in} E_1*E_2})=\unsure{\chi|\gamma}(u_1\otimes\chi(x_2}\\
&\rho(u_1\otimes x_2)=\chi u_1\otimes\bigl(\chi\deleted{(}x_2\deleted{)}-\psi_2(x_2)T_2\bigr)
\end{aligned}
\]
Mais
\[
\mu_{\deleted{\unsure{E_2|E_1,E_2}}}(u_1\otimes x_2)\overset{\mathrm{def}}{=}u_1\otimes
\underbrace{(\chi\deleted{\unsure{x}}x_2-\psi_2(x_2)T_2)}_{\unsure{\in}}
\]
% lines 4–6, restriction and diagonal note at left
\unsure{On|il faut} prendra donc $\sigma=\chi\,\mathrm{id}_{\deleted{E}L_1\otimes L_2}^{*}$ pour avoir commutativité du côté $L_1\otimes E_2\subset E_1*E_2$, et de $=$ pour $E_1\otimes L_2$. On gagne !
% slanted
\margin{* C'est ce qu'il fallait prendre \unsure{a|à} priori, si on voulait que $\pi_T$ soit une $\chi$-\unsure{trivialisation|trivialis.}}
% lines 7–13, regular scalar
Si $\chi$ est regulier sur \deleted{$\unsure{L_1,L_2}$} $L_1\otimes L_2$ (p.ex. $\chi$ régulier \unsure{sur|dans} $O$ et $L_1,L_2$ plats sur $O$, donc $L_1\otimes L_2$ plat aussi...) alors on peut définir $T$ par
\[
\begin{aligned}
\chi T&=\deleted{\unsure{\rho}}\mu'(T_1\otimes T_2)\overset{\mathrm{déf}}{=}T_1\deleted{\otimes}\unsure{*|\wedge}T_2\\
\text{i.e.}\qquad T&=\frac{1}{\chi}T_1\deleted{\otimes}*T_2\qquad(!)
\end{aligned}
\]
% lines 14–17, projection and its rearrangement
\deleted{Considérons} En effet considérons \deleted{$\unsure{\pi_T}$}
\[
\begin{aligned}
\pi_T:\ E&\longrightarrow L_1\otimes L_2\\
x&\longmapsto\chi x-T\psi(x)\qquad\text{d'où}\qquad T\psi(x)=
\deleted{\unsure{\chi\pi_T(x)}}\inserted{\chi x-\pi_T(x)}
\end{aligned}
\]
% lines 18–21, heavily struck calculation at left and continuation at right
\deleted{\unsure{puis|posons} faisons $x=T$, on trouve, comme $\psi(T)=\chi$}
\[
\begin{aligned}
&\deleted{\chi\,\pi_T=\unsure{\chi} \mathbin{\mapsto}\chi^2x-\chi T\psi(x)}\\
&\deleted{\chi T=\chi T\unsure{.}}
\end{aligned}
\]
Faisons $x=T_1*T_2$, d'où
\[
\begin{aligned}
\psi(x)&=\deleted{\unsure{\chi}}\psi(T_1)\cdot\psi(T_2)=\chi^2\\
\text{d'où}\qquad\chi^2T&=\chi\,T_1*T_2-
\underbrace{\pi_T(T_1*T_2)}_{\pi_{T_1}(\underset{\shortparallel\,0}{T_1})\otimes\pi_{T_{\unsure{1|2}}}(\underset{\shortparallel\,0}{T_2})}
\end{aligned}
\]
% lines 22–25, diagonally struck lower paragraph and display
\deleted{Soient \unsure{$x_1,x_2$} $x_1\in\GammaP_1$, $x_2\in\GammaP_2$, tels q. $\psi_1(x_1)=1$, $\psi_2(x_2)=1$ prenons $x=x_1*x_2$, on aura}
\[
\deleted{T=\chi x_1*x_2-\underbrace{\pi_T(x_1*x_2)}_{(\chi x_1-T_1)(\chi x_2-T_2)}}
\]
% archive sheet number: 70
